\section{The literal matched determinant coefficient}
\label{sec:literal-coefficient}

Fix once and for all a nonzero integer shift \(h_0\).  We work on
the high-\(\beta\) TPC-32 matched packet
\begin{equation}
 Q_X=X^{267/400+o(1)},\qquad
 J_X=X^{133/400+o(1)}.
\label{eq:high-beta-scales}
\end{equation}
In particular \(J_X\ll Q_X\).  Put
\begin{equation}
 C_X=\lfloor J_X\rfloor,\qquad
 N_{0,X}=J_XQ_X^2.
\label{eq:matched-scales}
\end{equation}
Hence \(C_X\le J_X\) and \(C_X\ll Q_X\), so the imported TPC-32
support, singleton, and energy estimates apply in their stated
range.
All fixed support and smoothness data are suppressed from the
notation; implied constants may depend on them and on \(h_0\), but
not on \(X\).

The signed normalized-determinant coefficient is
\begin{equation}
 A_X(n):=A_{C_X}(n),\qquad n\in\Z,
\label{eq:literal-A}
\end{equation}
where \(A_{C_X}\) is formed only after the three matched physical
channels and every native contribution in the determinant fiber have
been summed.  Define
\begin{equation}
\begin{aligned}
 S_X&=\#\supp A_X,\\
 L_X&=\|A_X\|_\infty,\\
 M_X&=\|A_X\|_1=\sum_n|A_X(n)|,\\
 D_X&=\|A_X\|_2^2=\sum_n|A_X(n)|^2.
\end{aligned}
\label{eq:SLMD}
\end{equation}
When \(A_X=0\), all four quantities are zero.  Every division by
\(S_X\) below is restricted to the nonzero case.

\begin{proposition}[Verified TPC-32 envelope]
\label{prop:tpc32-envelope}
For the literal coefficient in \eqref{eq:literal-A},
\begin{align}
 S_X&\ll Q_X,\label{eq:support-upper}\\
 M_X&\ll X^{o(1)}N_{0,X},\label{eq:mass-upper}\\
 L_X&\ll
 X^{o(1)}\frac{N_{0,X}}{Q_X},\label{eq:atom-upper}\\
 D_X&\ll
 X^{o(1)}\frac{N_{0,X}^2}{Q_X}.
\label{eq:energy-upper}
\end{align}
\end{proposition}

\begin{proof}
TPC-32 proves that the determinant support is contained in a fixed
multiple of
\([ -Q_X,Q_X]\cap(\Z\setminus\{0\})\), giving
\eqref{eq:support-upper}.  It also constructs a post-fiber absolute
majorant \(A_{C_X}^+\) with
\[
 \sum_n A_{C_X}^+(n)\ll X^{o(1)}N_{0,X}
\]
and
\[
 \max_n A_{C_X}^+(n)
 \ll X^{o(1)}Q_X(C_X+J_X).
\]
Since \(C_X\le J_X\) and
\(N_{0,X}/Q_X=J_XQ_X\), the second estimate is
\[
 \max_n A_{C_X}^+(n)
 \ll X^{o(1)}\frac{N_{0,X}}{Q_X}.
\]
The literal inequalities
\[
 |A_X(n)|\le A_{C_X}^+(n)
\]
give \eqref{eq:mass-upper} and \eqref{eq:atom-upper}.  Finally,
\[
 D_X\le L_XM_X
\]
gives \eqref{eq:energy-upper}; this is also the energy estimate stated
in TPC-32 \citep{WangTPC32}.
\end{proof}

\begin{remark}[The atom bound is already normalized]
\label{rem:no-renormalization}
The bound \eqref{eq:atom-upper} is used only in the normalization in
which TPC-32 proves it for the complete matched coefficient.  It is
not divided by a number of channels, triplets, fibers, or
components, and no diagonal or componentwise normalization is
inserted.  Any alternative normalization would require a new literal
comparison theorem.
\end{remark}

\begin{remark}[Actual versus available support]
The number \(S_X\) is the exact nonzero support after determinant
aggregation.  The interval of length \(O(Q_X)\) supplies only the
upper bound \eqref{eq:support-upper}.  It does not say that
\(S_X\asymp Q_X\); a lower support bound will instead follow
conditionally from energy survival in \cref{cor:support-forced}.
\end{remark}
